% GENERATED FILE. Edit canonical lessons, then run npm run generate:books.
% canonical-source-hash: fnv1a64:402ae7682140807e
% canonical-lessons: GE-C482-verstauchen, GE-C482-prellung, GE-C482-verband, GE-C482-desinfizieren, GE-C482-narbe

\chapter{Kleine Verletzungen}
\label{ch:ge-kleine-verletzungen}
\hlchaptermodality{482}

\begin{chapteropening}
\textbf{By the end of this chapter:} \emph{I can talk about small injuries and first aid.}

Complete the last lesson of chapter 482: I can talk about small injuries and first aid.
\end{chapteropening}

\section[verstauchen]{verstauchen — to sprain}

\label{lesson:GE-C482-verstauchen}



\begin{quote}
Before the new one: say the German for to scratch, then the German for to put cream on.
\end{quote}

\subsection*{You'll want to know: verstauchen}
\textbf{verstauchen} — \textquotedblleft{}to sprain\textquotedblright{}.

\begin{quote}
Beim Fußball habe ich mir den Fuß verstaucht.
\end{quote}

\begin{quote}
\emph{I sprained my foot playing football.}
\end{quote}

With a dative of the person: \emph{sich den Fuß verstauchen}. Inseparable: \emph{hat verstaucht}.

\begin{cousinweb}
\emph{verstauchen} is built on \emph{ver-} and \emph{stauchen} \textquotedblleft{}to compress, to jolt\textquotedblright{} (EtymDB): a joint pushed out of shape.
\end{cousinweb}

\begin{grammarlens}[title={What you've built}]
One more word for telling, explaining and describing at B1.
\end{grammarlens}

\subsection*{Your turn}
\begin{itemize}
\raggedright
  \item \emph{Say it:} \emph{verstauchen}
  \item \emph{Say it:} \emph{verstauchen}, once more
  \item \emph{Say it:} a sentence of your own with \emph{verstauchen}
  \item \emph{Recall:} say the German for to scratch, then the German for to put cream on, then say \emph{verstauchen} again
\end{itemize}

\begin{tcolorbox}[breakable,colback=teal!4,colframe=teal!35!black,title={Before you move on}]
What does verstauchen mean? (\textquotedblleft{}To sprain\textquotedblright{}.) Say it once more.
\end{tcolorbox}
\section[die Prellung]{die Prellung — bruise}

\label{lesson:GE-C482-prellung}



\begin{quote}
Before the new one: say the German for to put cream on, then the German for to sprain.
\end{quote}

\subsection*{You'll want to know: die Prellung}
\textbf{die Prellung} — \textquotedblleft{}bruise\textquotedblright{}.

\begin{quote}
Nach dem Sturz hatte er eine Prellung am Bein.
\end{quote}

\begin{quote}
\emph{After the fall he had a bruise on his leg.}
\end{quote}

Plural: \emph{die Prellungen}. In speech also \emph{der blaue Fleck}.

\begin{cousinweb}
\emph{Prellung} is built on \emph{prellen} \textquotedblleft{}to strike, to bump\textquotedblright{} (EtymDB) with \emph{-ung}; the older origin of \emph{prellen} is not settled.
\end{cousinweb}

\begin{grammarlens}[title={What you've built}]
One more word for telling, explaining and describing at B1.
\end{grammarlens}

\subsection*{Your turn}
\begin{itemize}
\raggedright
  \item \emph{Say it:} \emph{die Prellung}
  \item \emph{Say it:} \emph{die Prellung}, once more
  \item \emph{Say it:} a sentence of your own with \emph{die Prellung}
  \item \emph{Recall:} say the German for to put cream on, then the German for to sprain, then say \emph{die Prellung} again
\end{itemize}

\begin{tcolorbox}[breakable,colback=teal!4,colframe=teal!35!black,title={Before you move on}]
What does die Prellung mean? (\textquotedblleft{}Bruise\textquotedblright{}.) Say it once more.
\end{tcolorbox}
\section[der Verband]{der Verband — bandage, dressing; association}

\label{lesson:GE-C482-verband}



\begin{quote}
Before the new one: say the German for to sprain, then the German for bruise.
\end{quote}

\subsection*{You'll want to know: der Verband}
\textbf{der Verband} — \textquotedblleft{}bandage, dressing; association\textquotedblright{}.

\begin{quote}
Die Krankenschwester legt einen Verband an.
\end{quote}

\begin{quote}
\emph{The nurse puts on a bandage.}
\end{quote}

Plural: \emph{die Verbände}. Also an association or federation: \emph{der Fußballverband}.

\begin{cousinweb}
\emph{Verband} is built on \emph{verbinden} \textquotedblleft{}to bind up, to bandage\textquotedblright{}, from \emph{binden} \textquotedblleft{}to bind\textquotedblright{}, inherited from Old High German \emph{bintan}, from Proto-Germanic \emph{bindaną} (EtymDB).
\end{cousinweb}

\begin{grammarlens}[title={What you've built}]
One more word for telling, explaining and describing at B1.
\end{grammarlens}

\subsection*{Your turn}
\begin{itemize}
\raggedright
  \item \emph{Say it:} \emph{der Verband}
  \item \emph{Say it:} \emph{der Verband}, once more
  \item \emph{Say it:} a sentence of your own with \emph{der Verband}
  \item \emph{Recall:} say the German for to sprain, then the German for bruise, then say \emph{der Verband} again
\end{itemize}

\begin{tcolorbox}[breakable,colback=teal!4,colframe=teal!35!black,title={Before you move on}]
What does der Verband mean? (\textquotedblleft{}Bandage, dressing; association\textquotedblright{}.) Say it once more.
\end{tcolorbox}
\section[desinfizieren]{desinfizieren — to disinfect}

\label{lesson:GE-C482-desinfizieren}



\begin{quote}
Before the new one: say the German for bruise, then the German for bandage.
\end{quote}

\subsection*{You'll want to know: desinfizieren}
\textbf{desinfizieren} — \textquotedblleft{}to disinfect\textquotedblright{}.

\begin{quote}
Der Arzt desinfiziert die Wunde.
\end{quote}

\begin{quote}
\emph{The doctor disinfects the wound.}
\end{quote}

Regular, Perfekt without \emph{ge-}: \emph{hat desinfiziert}.

\begin{cousinweb}
\emph{desinfizieren} is built on \emph{des-} \textquotedblleft{}un-\textquotedblright{} and \emph{infizieren} \textquotedblleft{}to infect\textquotedblright{}, a learned borrowing from Latin \emph{inficere} \textquotedblleft{}to stain, to taint\textquotedblright{} (Lewis \& Short).
\end{cousinweb}

\begin{grammarlens}[title={What you've built}]
One more word for telling, explaining and describing at B1.
\end{grammarlens}

\subsection*{Your turn}
\begin{itemize}
\raggedright
  \item \emph{Say it:} \emph{desinfizieren}
  \item \emph{Say it:} \emph{desinfizieren}, once more
  \item \emph{Say it:} a sentence of your own with \emph{desinfizieren}
  \item \emph{Recall:} say the German for bruise, then the German for bandage, then say \emph{desinfizieren} again
\end{itemize}

\begin{tcolorbox}[breakable,colback=teal!4,colframe=teal!35!black,title={Before you move on}]
What does desinfizieren mean? (\textquotedblleft{}To disinfect\textquotedblright{}.) Say it once more.
\end{tcolorbox}
\section[die Narbe]{die Narbe — scar}

\label{lesson:GE-C482-narbe}



\begin{quote}
Before the new one: say the German for bandage, then the German for to disinfect.
\end{quote}

\subsection*{You'll want to know: die Narbe}
\textbf{die Narbe} — \textquotedblleft{}scar\textquotedblright{}.

\begin{quote}
Seit dem Unfall hat er eine Narbe am Knie.
\end{quote}

\begin{quote}
\emph{Since the accident he has had a scar on his knee.}
\end{quote}

Plural: \emph{die Narben}.

\begin{cousinweb}
\emph{Narbe} goes back to Middle High German \emph{narwe}; its older origin is not settled, and EtymDB lists the German word without a source.
\end{cousinweb}

\begin{grammarlens}[title={What you've built}]
One more word for telling, explaining and describing at B1.
\end{grammarlens}

\subsection*{Your turn}
\begin{itemize}
\raggedright
  \item \emph{Say it:} \emph{die Narbe}
  \item \emph{Say it:} \emph{die Narbe}, once more
  \item \emph{Say it:} a sentence of your own with \emph{die Narbe}
  \item \emph{Recall:} say the German for bandage, then the German for to disinfect, then say \emph{die Narbe} again
\end{itemize}

\begin{tcolorbox}[breakable,colback=teal!4,colframe=teal!35!black,title={Before you move on}]
What does die Narbe mean? (\textquotedblleft{}Scar\textquotedblright{}.) Say it once more.
\end{tcolorbox}
